\SourceLeaf{333}

\[
\sin(-1)\sin(2+3)\sin(2+5) = \sin(-2)\sin(1+4)\sin(1+6),
\]

where the quantities

\[
-1, \qquad 2+3, \qquad 2+5; \qquad -2, \qquad 1+4, \qquad 1+6
\]

satisfy the same conditions as those expressed for the quantities

\[
+1, \qquad 3, \qquad 5; \qquad +2, \qquad 4, \qquad 6
\]

by equations (1), (2), so that all the conclusions also carry over.

\begin{center}
[6.]

[\emph{Handbuch}, pp.\ 278 to 281.]

\textbf{FURTHER ON ASSESSING THE PHYSICAL POSSIBILITY OF THE DATA IN POTHENOT'S PROBLEM.}
\end{center}

$a$, $b$, $c$, $A$, $B$, $C$, $\alpha$, $\beta$, $\gamma$, $\mathfrak{A}$, $\mathfrak{B}$, $\mathfrak{C}$ have the same meaning as above. Namely,

\begin{center}\begin{tabular}{@{}lp{9.8cm}@{}}
$a,b,c$ & sides;\\
$A,B,C$ & angles of the triangle formed by the given points (1), (2), (3), with radius $=r$;\\
$\alpha,\beta,\gamma,\rho$ & distances of the sought point (0) from $A,B,C$ and the centre of that circle;\\
$\mathfrak{A}$ & angle from $02$ to $03$, measured in the same way as $A$ from $12$ to $13$ (i.\,e. both from right to left or both the other way round), and similarly $\mathfrak{B},\mathfrak{C}$.
\end{tabular}\end{center}

Then, as already noted on p.\ 270 [of the notebook], we have

\[
\frac{\sin(\mathfrak{A}-A)}{a\alpha} = \frac{\sin(\mathfrak{B}-B)}{b\beta} = \frac{\sin(\mathfrak{C}-C)}{c\gamma} = \frac{rr-\rho\rho}{2\alpha\beta\gamma r},
\]

from which it follows that, according as (0) lies inside or outside that circle,

\[
\sin(\mathfrak{A}-A), \qquad \sin(\mathfrak{B}-B), \qquad \sin(\mathfrak{C}-C)
\]

all become positive or negative.
